% CLAUDE REPORT::WRBRIDGE-0001::2026-11-15T01:51-08:00
\documentclass[11pt]{article}
\usepackage{amsmath,amssymb,amsthm}
\usepackage[margin=1in]{geometry}
\newtheorem{theorem}{Theorem}
\newtheorem{corollary}{Corollary}
\title{WRBRIDGE-0001: Difficulty conservation for the worst-case derivation bridge}
\author{Claude scheduled routine (supporting agent; not an LSPV/TMR owner)}
\date{2026-11-15 01:51 PST}
\begin{document}
\maketitle

\paragraph{Status.} PROVED/EXACT (elementary): Theorem~1 and Corollary~1.
KNOWN THEOREM: Theorem~2 (Cook 1975 translation), cited literature.
COMPUTATIONAL: exact minimum tree-like Resolution sizes for three tiny UNSAT instances.
The bridge $\mathrm{SAT}\in\mathrm P\Rightarrow\overline W_R(N)\le\mathrm{poly}(N)$ remains OPEN for every fixed $R$ for which a superpolynomial lower bound is known.
\textbf{P versus NP remains OPEN.}

\section{Target}
For a fixed sound and complete refutation system $R$ with a polynomial-time checker,
\[
\overline W_R(N)=\max_{F\in\mathrm{UNSAT},\ |\langle F\rangle|\le N}\ \min_{\pi:F\vdash_R\bot}\operatorname{cost}_R(\pi).
\]
Taking the minimum over all derivations is the standard Cook--Reckhow proof-length measure. Not charging for finding $\pi$ is legitimate for a lower bound, and the corpus already uses it (MAIN 317, MAIN 344).
The \emph{bridge} is $\mathrm{Bridge}(R)$: $\mathrm{SAT}\in\mathrm P\Rightarrow\overline W_R(N)\le\mathrm{poly}(N)$.

\section{Difficulty conservation}
\begin{theorem}[PROVED/EXACT]
For every fixed $R$, $\mathrm{Bridge}(R)$ is logically equivalent to $\mathrm{SAT}\notin\mathrm P\ \vee\ \overline W_R\le\mathrm{poly}$.
So if $\overline W_R(N)\ge N^{\omega(1)}$ is proved unconditionally, then $\mathrm{Bridge}(R)\iff \mathrm P\ne\mathrm{NP}$.
\end{theorem}
\begin{proof}
The first statement is the definition of material implication. If $\overline W_R$ is not polynomially bounded, the second disjunct is false, so the equivalence reduces to $\mathrm{SAT}\notin\mathrm P$.
\end{proof}
\begin{corollary}[PROVED/EXACT]
Superpolynomial lower bounds on $\overline W_R$ are KNOWN THEOREMS for these systems: tree-like and DAG Resolution (Haken 1985; tree-like pigeonhole $2^{\Omega(n\log n)}$, Iwama--Miyazaki 1999), Cutting Planes (Pudl\'ak 1997), bounded-depth Frege (Ajtai; Kraj\'i\v{c}ek--Pudl\'ak--Woods; Pitassi--Beame--Impagliazzo) and Polynomial Calculus (Razborov 1998).
For each of them, proving the bridge is exactly as hard as proving $\mathrm P\ne\mathrm{NP}$.
So the plan ``lower bound for $R$ plus bridge for $R$'' cannot be split into two strictly easier halves for any system whose lower-bound half is currently within reach.
All of the difficulty sits in the bridge.
\end{corollary}
\paragraph{Consequence for $R_{\mathrm{mix341}}$ (CANDIDATE classification).}
Suppose $R_{\mathrm{mix341}}$ p-simulates Resolution and no superpolynomial lower bound for it is in the corpus.
Then the equivalence in Theorem~1 is not yet triggered, but a lower bound for $R_{\mathrm{mix341}}$ would trigger it.
This is a structural reason to expect that the B3 transfer (MAIN 317: B3 $\iff$ length simulation of the decider-induced system) is the load-bearing arrow.

\paragraph{The free-sequence clarification has a direction.}
Results showing that automating Resolution, Cutting Planes or algebraic systems is NP-hard (Atserias--M\"uller 2019; G\"o\"os et al.\ 2020; de~Rezende et al.\ 2021) concern the direction from short proofs to algorithms.
The bridge needs the opposite direction, from a polynomial-time algorithm to short $R$-proofs, which is p-simulation of the decider-induced system $R_A$ by $R$.
Search hardness therefore neither helps nor hurts the bridge, which is consistent with the user's clarification.

\section{The only known route to a bridge: feasible provability}
\begin{theorem}[KNOWN THEOREM; Cook 1975 PV--EF translation]
Let $A$ be a polynomial-time SAT decider such that $\mathrm{PV}\vdash\forall F\,\forall a\,(A(F)=0\rightarrow \neg\mathrm{Sat}(F,a))$.
Then Extended Frege has polynomial-size refutations of every UNSAT $F$, so $\overline W_{\mathrm{EF}}(N)\le\mathrm{poly}(N)$.
\end{theorem}
\emph{Sketch.} The PV theorem translates into polynomial-size EF proofs of $\|A(x)=0\to\neg\mathrm{Sat}(x,y)\|_n$. Substitute the bits of $F$ and verify the run $A(F)=0$ inside EF in polynomial size.

So $\mathrm{Bridge}_{\mathrm{PV}}(\mathrm{EF})$ (``a PV-provably correct polynomial-time SAT algorithm implies EF is p-bounded'') is the weakest bridge currently known to be provable.
Its contrapositive with an EF lower bound gives only ``no PV-provably correct polynomial-time SAT algorithm''. That is an unprovability result, not $\mathrm P\ne\mathrm{NP}$.
Pich--Santhanam, \emph{Towards P$\ne$NP from Extended Frege lower bounds} (arXiv:2312.08163), is the state of the art on removing the PV hypothesis.
They show that EF not being p-bounded implies $\mathrm P\ne\mathrm{NP}$ under extra $S^1_2$-provability hypotheses, for example that $S^1_2$ proves a polynomial-time function generates anticheckers for SAT, or that it witnesses errors of small circuits for SAT.
Exact object correspondence: their ``EF not p-bounded'' is our $\overline W_{\mathrm{EF}}\ne\mathrm{poly}$ with $\operatorname{cost}=$ proof length. Their hypotheses are exactly the missing transfer theorem.
Classification: DIRECTLY TRANSFERABLE THEOREM (conditional).

\paragraph{Barriers for the lower-bound half at EF strength (KNOWN).}
EF has no feasible interpolation if RSA or Diffie--Hellman is secure (Kraj\'i\v{c}ek--Pudl\'ak 1998; Bonet--Pitassi--Raz 2000).
Pich--Santhanam (FOCS 2019) prove that proof systems cannot efficiently prove circuit lower bounds for random truth tables, which is a natural-proofs analogue.
Arteche--Atserias--de~Rezende--Khaniki 2025 (arXiv:2506.16956) give PV-formalized Resolution non-automatability and show that the proof analysis problem for EF is NP-complete.
This confirms the feasible-provability lens is active current work.
PQ-ledger item 5 (worst-case/average-case): FAILED TRANSFER here, because the cryptographic assumptions are inadmissible in an unconditional proof and they only block feasible interpolation.

\section{Computational fixtures (COMPUTATIONAL)}
The exact minimum number of leaves of a tree-like Resolution refutation equals the minimum number of leaves of a decision tree for the falsified-clause search problem.
Code: \texttt{experiments/wr\_bridge/2026-11-15\_treeres\_min.py}. Raw data: \texttt{results/raw/wr\_bridge/2026-11-15\_treeres\_min.json}.
Each instance is one stage of under 3\,s, and the two 6-variable instances were cross-checked by a solver without pruning.
\begin{center}
\begin{tabular}{lrrrr}
Instance & vars & clauses & min leaves & min tree size\\\hline
$\mathrm{PHP}^3_2$ & 6 & 9 & 11 & 21\\
Tseitin $K_4$ (odd charge) & 6 & 16 & 28 & 55\\
$\mathrm{PHP}^4_3$ & 12 & 22 & 43 & 85
\end{tabular}
\end{center}
These are certified shortest-derivation fixtures for the Builder and Red Team roles. They say nothing about asymptotics.

\section{Role assignments derived this run}
\begin{itemize}
\item \textbf{General SAT Lower-Bound Lead:} target $\mathrm{Bridge}_{\mathrm{PV}}$ or one Pich--Santhanam hypothesis for $R\supseteq\mathrm{EF}$. Theorem~1 rules out treating the bridge as weaker than P$\ne$NP for Resolution-strength systems.
\item \textbf{SAT-Circuit Reduction Builder:} decide whether $R_{\mathrm{mix341}}$ p-simulates EF. If it does not, Theorem~2 gives no bridge for it.
\item \textbf{Hypergraph \& Proof Certificates:} hard families already have lower bounds for Resolution, CP and PC. A new family is useful only if it is a candidate for EF hardness, and the 2019 barrier applies there.
\item \textbf{Time-Space Resource Auditor:} cost in Theorem~2 is EF length. Audit the polynomial overhead of the PV-to-EF translation against any fixed exponent ledger.
\item \textbf{Independent Proof Red Team:} test every proposed bridge for $R$ against Theorem~1. If $R$ has a known lower bound, the proposal is a disguised proof of P$\ne$NP and must be audited as such.
\end{itemize}

\section*{CORPUS REVIEWED}
\texttt{AGENTS.md}; the ownership map of root \texttt{main.tex}; the git log through TMR-3/040 (c6d3b05d) and GROKPROOF 2026-11-15 08:20; MAIN 317 (Cook--Reckhow B3 transfer); MAIN 344 (p-optimal collapse); LSPV-1/337 (finite reflection bridge); the prior Claude world ledger (row L13).
\section*{DUPLICATION / CONTRADICTION AUDIT}
MAIN 317 already proves B3 $\iff$ length simulation of the decider-induced system. Theorem~1 is its disjunctive and contrapositive form; the increment is the explicit consequence for systems with known lower bounds.
MAIN 344 already covers the p-optimal equivalence under SAT$\in$P.
Theorem~2 and the Pich--Santhanam mapping were not found by a repository search for ``Krajicek-Pudlak'', ``p-bounded'' and ``Bridge(R)''. Those hits were ER/reflection reports, not the PV-provability route.
No contradiction was found.
\section*{WHAT CHANGED RELATIVE TO ROOT MAIN}
Nothing. Root \texttt{main.tex} has no Claude-owned region, and the protocol forbids editing LSPV/TMR regions. This report and \texttt{research/agents/claude/LIVING\_LOG.md} are the living record.
\section*{ADOPTED / REJECTED / RETIRED PRIOR COMPONENTS}
Adopted: the B3 framing from MAIN 317. Rejected: treating $\mathrm{Bridge}(\mathrm{Res})$ as a weaker target than P$\ne$NP.
\section*{STALE OR OVERSTATED CLAIMS FOUND}
None found in the files reviewed.

\end{document}
% END_OF_VERIFIED_REPORT::CLAUDE::WRBRIDGE-0001
